\chapter{D-Dimer}

\section*{What Is This Test?}
D-dimer is a fibrin degradation product (FDP) that consists of two cross-linked D-domains of fibrin generated when the fibrinolytic enzyme plasmin cleaves cross-linked fibrin. Unlike other fibrin degradation products, D-dimer is specific for the degradation of cross-linked fibrin (which requires prior thrombin generation, fibrinogen-to-fibrin conversion, and factor XIIIa-mediated cross-linking), distinguishing it from fibrinogen degradation products that can be produced without prior thrombus formation. The presence of D-dimer in the circulation therefore reflects both thrombin generation (coagulation activation) and plasmin generation (fibrinolysis). D-dimer is the most clinically useful biomarker for venous thromboembolism (VTE), including deep vein thrombosis (DVT) and pulmonary embolism (PE). Its primary clinical role is as a rule-out test: a negative D-dimer result in a patient with low-to-intermediate pre-test probability (assessed by validated clinical prediction rules such as the Wells score) effectively excludes acute VTE without the need for further imaging. The D-dimer is also used in the diagnosis and monitoring of disseminated intravascular coagulation (DIC) and as a prognostic marker in sepsis, trauma, and other conditions with activation of coagulation and fibrinolysis. D-dimer levels increase with age, pregnancy, surgery, cancer, and inflammation, so specificity for VTE is lower in these populations, and age-adjusted or clinical-context-adjusted cutoffs have been developed to improve diagnostic performance.

\section*{Alternative Names}
Fibrin D-dimer; D-dimer Test; Fibrin Degradation Product; D-dimer Assay; Cross-linked Fibrin Degradation Product

\section*{Principle}
D-dimer is measured by enzyme-linked immunosorbent assay (ELISA), latex-enhanced immunoturbidimetric assay, or whole-blood agglutination methods, using monoclonal antibodies specific for the D-dimer epitope (which is formed by the cross-linking of D-domains by factor XIIIa and is not present on fibrinogen or non-cross-linked fibrin degradation products). Quantitative automated latex-enhanced immunoturbidimetric assays (e.g., HemosIL D-dimer on IL ACL systems, Innovance D-dimer on Sysmex, STA-Liatest D-dimer on Stago) are the most widely used in clinical laboratories. In these assays, patient plasma is mixed with latex microparticles coated with monoclonal anti-D-dimer antibodies. If D-dimer is present, it cross-links the particles, causing agglutination. The increase in turbidity is measured spectrophotometrically and is proportional to D-dimer concentration. Results are reported in fibrinogen equivalent units (FEU) or D-dimer units (DDU), with 1 FEU approximately equal to 2 DDU (since D-dimer has approximately half the molecular weight of fibrinogen, 1 ng/mL D-dimer is 2 ng/mL FEU). Different manufacturers use different monoclonal antibodies recognizing different D-dimer epitopes, different calibration standards, and different reporting units, making results from different assays non-interchangeable. The conventional cutoff for VTE exclusion is typically 500 ng/mL FEU (or 250 ng/mL DDU). Age-adjusted cutoffs (patient age $\times$ 10 ng/mL for patients >50 years in FEU) have been validated to improve specificity while maintaining high sensitivity.

\section*{History}
D-dimer was first identified as a specific degradation product of cross-linked fibrin in the 1970s, distinguishing it from other fibrin/fibrinogen degradation products that lack the cross-link structure. Monoclonal antibodies specific for the D-dimer epitope were developed in the 1980s, enabling the development of specific, reproducible D-dimer assays. The first clinical studies establishing D-dimer as a rule-out test for VTE were published in the late 1980s and early 1990s. In 1995, the landmark study by Wells et al. combined D-dimer testing with clinical probability assessment (the Wells score), establishing the modern diagnostic algorithm for VTE. The ADJUST-PE trial (2014) validated the age-adjusted D-dimer cutoff (age $\times$ 10 in patients >50 years), which has since been incorporated into clinical guidelines. The YEARS algorithm (2017) further simplified the diagnostic approach by combining D-dimer thresholds with a streamlined clinical decision rule.

\section*{Why This Test Is Ordered}
\begin{itemize}
  \item Exclusion of venous thromboembolism (DVT and PE) in patients with low-to-intermediate pre-test probability --- a negative D-dimer (<500 ng/mL FEU or below age-adjusted cutoff) effectively excludes acute VTE without imaging
  \item Diagnostic workup of suspected pulmonary embolism in the emergency department, combined with clinical decision rules (Wells score, YEARS criteria, Geneva score)
  \item Diagnostic workup of suspected deep vein thrombosis (lower extremity, upper extremity)
  \item Diagnosis and monitoring of disseminated intravascular coagulation (DIC) --- marked D-dimer elevation is a key criterion in the ISTH DIC scoring system
  \item Assessment of thrombotic risk and monitoring of DIC in sepsis, trauma, major surgery, and obstetric emergencies (HELLP syndrome, amniotic fluid embolism, placental abruption)
  \item Evaluation of suspected acute aortic dissection --- elevated D-dimer at a very high cutoff (5,000 ng/mL) has been proposed as a rule-in test
  \item Prognostic assessment in COVID-19-associated coagulopathy
  \item Monitoring of anticoagulation cessation --- D-dimer elevation after stopping anticoagulation identifies patients at higher risk of recurrent VTE
  \item Investigation of unexplained elevation of coagulation markers
\end{itemize}

\section*{Is This a Primary or Secondary Test}
D-dimer is a primary screening test for the exclusion of venous thromboembolism in patients with low-to-intermediate clinical probability and a key test in the diagnosis of DIC and monitoring of fibrinolytic states.

\section*{How This Test Is Performed}
A venous blood sample is collected in a light blue-top tube containing 3.2\% sodium citrate (9:1 blood-to-citrate ratio), filled exactly to the line. After gentle inversion, the sample is centrifuged at 1,500--2,500 $\times$ g for 15 minutes to obtain platelet-poor plasma. The plasma is aspirated and analyzed on an automated coagulation analyzer using a latex-enhanced immunoturbidimetric method. Results are generated within 10--20 minutes. Venipuncture takes less than 5 minutes. Point-of-care D-dimer assays using whole blood are also available.

\section*{What Preparation Is Needed}
No special preparation is required. Fasting is not necessary. The patient should disclose all medications, particularly anticoagulants (heparin, warfarin, DOACs), which reduce D-dimer levels and may cause a false-negative result in patients who have started anticoagulation before testing. Recent surgery, trauma, pregnancy, malignancy, and inflammatory conditions should be disclosed, as these elevate D-dimer levels and reduce the test's specificity.

\section*{Contraindications and Precautions}
No absolute contraindications. Critical considerations for correct interpretation: D-dimer is a highly sensitive but nonspecific biomarker --- numerous conditions cause D-dimer elevation (infection, inflammation, malignancy, pregnancy, post-operative state, trauma, liver disease, advanced age, renal failure, and hospitalization itself). The specificity of D-dimer for VTE decreases with age (hence age-adjusted cutoffs), pregnancy, cancer, and hospitalization (where D-dimer is elevated in the majority of patients, limiting its utility). D-dimer should NEVER be used alone to diagnose VTE --- it is a rule-out test, not a rule-in test. A positive D-dimer in a patient with low pre-test probability requires confirmatory imaging (compression ultrasound for DVT, CT pulmonary angiography or V/Q scan for PE). A negative D-dimer in a patient with HIGH pre-test probability for PE (Wells score >6) does NOT safely exclude PE, and imaging should be performed. Anticoagulant therapy (heparin, DOACs) started before D-dimer testing can reduce D-dimer, creating a false-negative result. Different D-dimer assays use different antibodies, calibrators, and units (FEU vs. DDU) and are NOT interchangeable; the appropriate cutoff for the specific assay must be used. A normal D-dimer is assay-specific. D-dimer levels have limited clinical utility in hospitalized patients, intensive care unit patients, and immediate post-operative patients (where the vast majority will have elevated D-dimer, and a negative result may still be useful for VTE exclusion). Samples should be processed within 4 hours of collection or plasma separated and frozen.

\section*{Risks}
Minimal risks of venipuncture: mild pain, bruising, hematoma, lightheadedness, and rarely infection or nerve injury.

\section*{What Happens After the Test}
Results are available within 30--60 minutes. A negative D-dimer in a patient with low or unlikely pre-test probability excludes VTE, and no further VTE-specific imaging is needed. A positive D-dimer in any patient requires confirmatory imaging. A markedly elevated D-dimer (>4--5 times the upper limit of normal) in a patient with compatible clinical presentation (chest pain, hypotension, pulse deficit) raises concern for acute aortic dissection, and emergent CT angiography of the chest/abdomen should be pursued. In DIC evaluation, D-dimer is markedly elevated (often >10,000 ng/mL FEU) and is interpreted with platelet count, PT, aPTT, and fibrinogen. In COVID-19, markedly elevated D-dimer (>3,000--5,000 ng/mL) is associated with severe disease and increased mortality. Serial D-dimer measurements in DIC or COVID-19 can monitor response to therapy.

\section*{Understanding Results}

\subsection*{Below Normal (Negative D-dimer: <500 ng/mL FEU)}
\begin{itemize}
  \item \textbf{Normal finding:} Indicates absence of significant thrombin generation and fibrinolysis. Effectively excludes acute VTE in patients with low-to-intermediate pre-test probability (negative predictive value >99\%).
  \item \textbf{On therapeutic anticoagulation:} Heparin, warfarin, or DOAC therapy reduce thrombin generation, potentially lowering D-dimer even in the presence of VTE (but anticoagulation should not be started before diagnostic testing when VTE is suspected).
  \item \textbf{Small, distal, non-occlusive thrombus:} Rare false-negative result; a small calf DVT may produce D-dimer below assay detection threshold.
  \item \textbf{Delayed testing:} D-dimer may decrease if testing is performed >5--7 days after symptom onset (fibrinolytic activity diminishes as thrombus organizes).
\end{itemize}

\subsection*{Above Normal (Positive D-dimer: >500 ng/mL FEU)}
\begin{itemize}
  \item \textbf{Venous thromboembolism (PE, DVT):} D-dimer is elevated in >95\% of acute VTE. The magnitude of elevation correlates with clot burden (massive PE has higher D-dimer than small subsegmental PE), but the specific level cannot reliably distinguish PE from DVT or determine thrombus extent.
  \item \textbf{Disseminated intravascular coagulation (DIC):} Markedly elevated D-dimer (often >10,000--20,000 ng/mL FEU) reflects widespread microvascular thrombosis and secondary fibrinolysis. Combined with thrombocytopenia and prolonged PT/aPTT.
  \item \textbf{Sepsis and systemic infection:} Activation of coagulation and fibrinolysis in severe sepsis; D-dimer correlates with severity and mortality.
  \item \textbf{Major surgery and trauma:} Post-operative D-dimer elevation peaks at 3--7 days and may remain elevated for weeks, limiting specificity for VTE.
  \item \textbf{Malignancy:} Cancer-associated activation of coagulation (Trousseau syndrome). D-dimer is often elevated in metastatic disease and is a poor prognostic factor.
  \item \textbf{Pregnancy and postpartum period:} Progressive D-dimer elevation throughout pregnancy, with normal values reaching 2--3 times non-pregnant levels in the third trimester. Further elevated in pre-eclampsia, HELLP syndrome, and obstetric DIC.
  \item \textbf{Acute aortic dissection:} D-dimer is markedly elevated; a very high cutoff (5,000 ng/mL) has been proposed as a rule-in test with high sensitivity for aortic dissection.
  \item \textbf{Cardiovascular disease:} Acute myocardial infarction, acute coronary syndrome, atrial fibrillation, heart failure, and peripheral arterial disease.
  \item \textbf{Stroke:} Acute ischemic stroke and hemorrhagic stroke. D-dimer correlates with infarct volume and outcome.
  \item \textbf{COVID-19:} Markedly elevated D-dimer, associated with disease severity, thrombosis risk, and mortality. Thresholds of 1,000--3,000 ng/mL used to guide escalation of anticoagulation in some protocols.
  \item \textbf{Liver disease:} Decreased clearance of fibrin degradation products by the liver.
  \item \textbf{Renal failure:} Decreased renal clearance of D-dimer fragments.
  \item \textbf{Advanced age:} D-dimer increases with age, necessitating the use of age-adjusted cutoffs (age $\times$ 10 ng/mL for patients >50 years, FEU) to maintain specificity.
  \item \textbf{Hematoma, ecchymosis, and soft tissue injury:} Fibrin formation at sites of hemorrhage or tissue injury.
  \item \textbf{Venous malformations and vascular aneurysms:} Activation of coagulation within abnormal vascular structures.
  \item \textbf{Thrombolytic therapy:} Massive D-dimer elevation following tPA or streptokinase administration for thrombolysis.
\end{itemize}

\section*{Normal Results}
\begin{itemize}
  \item \textbf{Adults (<50 years):} <500 ng/mL FEU (<0.5 mg/L FEU) or <250 ng/mL DDU --- reference range varies by specific assay
  \item \textbf{Adults (>50 years):} Age-adjusted cutoff = Age $\times$ 10 ng/mL FEU (e.g., for a 70-year-old patient: normal if <700 ng/mL FEU). Validated in the ADJUST-PE trial and incorporated into clinical guidelines (ACCP, ASH, ESC).
  \item \textbf{Pregnancy:} Progressive physiological elevation; first trimester <650 ng/mL, second trimester <1,200 ng/mL, third trimester <2,500 ng/mL (approximate, assay-dependent). A specific pregnancy-adapted D-dimer cutoff has not been universally adopted.
  \item \textbf{Children:} <500 ng/mL FEU (similar to adult non-pregnant range after the neonatal period)
  \item \textbf{Note:} D-dimer results from different manufacturers are NOT interchangeable. The specific assay used must be identified, and the appropriate cutoff for that assay must be applied.
\end{itemize}

\section*{Follow-up Tests}
\begin{itemize}
  \item \textbf{Compression ultrasound (CUS) of lower extremities:} First-line imaging for suspected DVT when D-dimer is positive
  \item \textbf{CT pulmonary angiography (CTPA):} First-line imaging for suspected PE when D-dimer is positive
  \item \textbf{Ventilation-perfusion (V/Q) scan:} Alternative to CTPA for PE when CT contrast is contraindicated (renal failure, severe contrast allergy) or to minimize radiation (young adults, pregnancy)
  \item \textbf{Wells score or Geneva score for pre-test probability of PE:} Must be calculated before D-dimer testing; high pre-test probability should proceed directly to imaging
  \item \textbf{Complete blood count with platelet count:} For DIC evaluation (thrombocytopenia)
  \item \textbf{PT/INR and aPTT:} For DIC evaluation (prolonged)
  \item \textbf{Fibrinogen (Clauss method):} For DIC evaluation (decreased)
  \item \textbf{Peripheral blood smear:} For DIC (schistocytes, microangiopathic hemolytic anemia)
  \item \textbf{CT angiography of chest/abdomen:} For suspected aortic dissection when D-dimer is markedly elevated with compatible presentation
  \item \textbf{Echocardiography:} For right ventricular strain assessment in PE, aortic dissection, or cardiac disease
  \item \textbf{Repeat D-dimer after anticoagulation cessation:} To assess risk of VTE recurrence
  \item \textbf{Septic workup:} Blood cultures, lactate, CRP, procalcitonin when elevated D-dimer in suspected sepsis
\end{itemize}

\section*{References}
\begin{enumerate}
  \item Tierney LM, Saint S, Drazen JM. CURRENT Medical Diagnosis \& Treatment. McGraw-Hill; 2023.
  \item Wallach J. Interpretation of Diagnostic Tests. 11th ed. Wolters Kluwer; 2022.
  \item Fischbach FT, Dunning MB. A Manual of Laboratory and Diagnostic Tests. 11th ed. Wolters Kluwer; 2023.
  \item Burtis CA, Ashwood ER, Bruns DE. Tietz Textbook of Clinical Chemistry and Molecular Diagnostics. 7th ed. Elsevier; 2023.
  \item Wells PS, Anderson DR, Rodger M, et al. Evaluation of D-dimer in the diagnosis of suspected deep-vein thrombosis. N Engl J Med. 2003;349(13):1227--1235.
  \item Righini M, Van Es J, Den Exter PL, et al. Age-adjusted D-dimer cutoff levels to rule out pulmonary embolism: the ADJUST-PE study. JAMA. 2014;311(11):1117--1124.
  \item van der Hulle T, Cheung WY, Kooij S, et al. Simplified diagnostic management of suspected pulmonary embolism (the YEARS study): a prospective, multicentre, cohort study. Lancet. 2017;390(10091):289--297.
  \item Lim W, Le Gal G, Bates SM, et al. ASH 2018 guidelines for management of venous thromboembolism: diagnosis of venous thromboembolism. Blood Adv. 2018;2(22):3226--3256.
  \item Adam SS, Key NS, Greenberg CS. D-dimer antigen: current concepts and future prospects. Blood. 2009;113(13):2878--2887.
\end{enumerate}

\clearpage
